\chapter{Analyse}

\section{Moteurs disponibles}

\logiciel{} peut s'appuyer sur plusieurs moteurs de calcul :

\begin{table}[H]
  \centering
  \begin{tabular}{lll}
    \toprule
    Moteur & Statut & Usage principal \\
    \midrule
    PyNite & Par défaut & Calcul statique linéaire courant \\
    OpenSeesPy & Optionnel & Calcul avancé et fonctionnalités spécifiques \\
    \bottomrule
  \end{tabular}
  \caption{Moteurs de calcul}
\end{table}

\section{Analyse statique linéaire}

L'analyse statique linéaire calcule les déplacements, réactions et efforts
internes pour les cas de charge et combinaisons disponibles.

Avant de lancer l'analyse, vérifier :

\begin{itemize}
  \item que le modèle contient au moins des nœuds et des éléments ;
  \item qu'au moins un appui est défini ;
  \item qu'il existe des charges ou des combinaisons ;
  \item que les matériaux et sections sont affectés correctement ;
  \item que le moteur choisi supporte les objets du modèle.
\end{itemize}

\section{Poids propre}

Le cas \fichier{Poids propre} est géré automatiquement. Il est calculé à partir
des sections, matériaux et densités disponibles.

\section{Messages de calcul}

La console inférieure affiche les informations importantes :

\begin{itemize}
  \item moteur utilisé ;
  \item cas calculés ;
  \item succès ou échec ;
  \item messages d'erreur éventuels.
\end{itemize}

\begin{warningbox}
Un calcul qui converge n'est pas automatiquement un calcul valide. Il faut
toujours vérifier la géométrie, les appuis, les charges, les unités et les
conventions de signe.
\end{warningbox}
